\section{Implementation}

\begin{frame}[shrink=7]{How the Repository Mirrors the Math}
    \begin{columns}[T,onlytextwidth]
        \begin{column}{0.56\textwidth}
            \begin{tikzpicture}[
    scale=0.88,
    transform shape,
    node distance=0.55cm and 0.8cm,
    file/.style={rounded corners=8pt, draw=black!75, line width=0.9pt, minimum width=3.8cm, minimum height=0.95cm, align=left, font=\scriptsize, inner sep=6pt},
    arrow/.style={-{Latex[length=3mm]}, line width=1.15pt}
]
    \node[file, fill=blue!10] (core) at (0,0) {\texttt{src/gmm\_missing.py}\\alternating estimator};
    \node[file, fill=green!12, below=of core] (exp1) {\texttt{notebooks/step\_by\_step\_experiment.ipynb}\\single-cycle explanation};
    \node[file, fill=orange!17, below=of exp1] (exp2) {\texttt{notebooks/GMM\_Nhat.ipynb}\\CC GENERAL reproduction};
    \node[file, fill=red!10, below=of exp2] (exp3) {\texttt{src/airplane\_crashes\_pipeline.py}\\real-data workflow};
    \node[file, fill=gray!10, right=1.6cm of exp2, minimum width=3.4cm, minimum height=4.8cm, align=center] (report) {report \& slides\\[0.2em]
    formulas\\figures\\artifacts\\discussion};

    \draw[arrow] (core) -- (report);
    \draw[arrow] (exp1.east) -- (report.west);
    \draw[arrow] (exp2.east) -- (report.west);
    \draw[arrow] (exp3.east) -- (report.west);
\end{tikzpicture}

        \end{column}
        \begin{column}{0.4\textwidth}
            \begin{beamercolorbox}[sep=0.8em, rounded=true, shadow=true]{block title}
                \textbf{Core file}
            \end{beamercolorbox}
            \vspace{0.15cm}
            \texttt{src/gmm\_missing.py}

            \vspace{0.25cm}
            \small
            \begin{beamercolorbox}[sep=0.65em, rounded=true, shadow=true]{block body}
                \textbf{Numerical safeguards}
                \begin{itemize}
                    \item mean fill only for initialization
                    \item KMeans initialization
                    \item covariance ridge regularization
                    \item log-sum-exp in E-step
                    \item fallback when precision blocks are singular
                \end{itemize}
            \end{beamercolorbox}
        \end{column}
    \end{columns}
\end{frame}


\begin{frame}[shrink=4]{Code-to-Equation Mapping}
    \scriptsize
    \renewcommand{\arraystretch}{1.45}
    \begin{tabular}{>{\bfseries}p{0.2\textwidth} p{0.33\textwidth} p{0.38\textwidth}}
        Equation block & Meaning & In code \\
        \toprule
        \(\gamma_{jk}\) & posterior responsibilities & \texttt{\_e\_step()} with log-domain Gaussian scores \\
        \(\alpha_k,\mu_k,\Sigma_k\) & standard GMM M-step & \texttt{\_m\_step\_params()} \\
        \(x_{j,m}^{\star}\) & precision-weighted update for missing entries & \texttt{\_m\_step\_data()} \\
        fit loop & alternating ascent until tolerance & \texttt{fit()} / \texttt{fit\_scale()} \\
        \bottomrule
    \end{tabular}

    \vspace{0.45cm}
    \begin{beamercolorbox}[sep=0.8em, rounded=true, shadow=true]{title}
        \large
        The implementation is not a black box.\\
        \textbf{The paper equations are visible almost line by line.}
    \end{beamercolorbox}
\end{frame}
